\documentclass[11pt,a4paper]{article}

\usepackage[utf8]{inputenc}
\usepackage[T1]{fontenc}
\usepackage{amsmath,amssymb}
\usepackage{graphicx}
\usepackage{booktabs}
\usepackage{siunitx}
\usepackage[margin=2.5cm]{geometry}
\usepackage[hidelinks]{hyperref}
\usepackage{float}
\usepackage{caption}
\usepackage{subcaption}

\sisetup{
  round-mode=places,
  round-precision=2,
  group-separator={,},
  group-minimum-digits=4,
}

\title{AzONetV2: UNET Inference Results \\ {\normalsize Statistical Analysis of Thickness Predictions}}
\author{}
\date{\today}

\begin{document}
\maketitle

\begin{abstract}
This report presents a comprehensive evaluation of the UNET-based thickness prediction model developed for the AzONetV2 project. Predictions were generated for three simulated data splits (960,000 train, 120,000 validation, and 120,000 test samples) and one external experimental dataset (145 thin-film samples at 145\,nm). Error metrics (MAE, MSE, RMSE, MAPE, $R^2$) and detailed distribution analyses of absolute percentage error (APE) and percentage error (PE) are reported for each split. The model achieves exceptional performance on simulated data (MAPE $\approx 0.61\%$, $R^2 = 0.9999$) with near-identical metrics across splits, indicating excellent generalization and no evidence of overfitting. On experimental 145\,nm samples, the model yields a MAPE of 7.38\% and $R^2 = 0.9385$.
\end{abstract}

\setcounter{tocdepth}{1}
\tableofcontents
\newpage

%=====================================================================
\section{Introduction}
%=====================================================================

The AzONetV2 project addresses the problem of estimating thin-film thickness from reflectance/transmission spectra using deep learning. A 1D UNET architecture was trained end-to-end to map normalized spectral inputs to thickness values. This report documents the final inference stage, where the trained model is evaluated on both simulated and experimental data, and a full battery of error metrics and statistical tests is computed.

The inference pipeline consists of two scripts located in \texttt{notebooks/7\_Final\_Inference/}:

\begin{itemize}
  \item \texttt{1\_figures\_MSE.py} — Loads pre-computed predictions (saved as \texttt{.npy} files by the \texttt{6\_results\_colab} stage), computes error metrics and distribution analyses, and generates publication-quality figures.
  \item \texttt{training\_plots.py} — Reads the training log (\texttt{training.log}) and plots the evolution of MSE, MAPE, and RMSE over 100 training epochs.
\end{itemize}

%=====================================================================
\section{Training History}
%=====================================================================

The UNET model was trained for 100 epochs with a fixed learning rate of $10^{-3}$. Table~\ref{tab:training_final} shows the final training and validation metrics, and Figure~\ref{fig:training} illustrates the full learning curves.

\begin{table}[H]
\centering
\caption{Final training metrics at epoch 100.}
\label{tab:training_final}
\begin{tabular}{lccc}
\toprule
\textbf{Metric} & \textbf{Train} & \textbf{Validation} \\
\midrule
MSE (loss)    & 660.20\,nm\textsuperscript{2} & 260.79\,nm\textsuperscript{2} \\
MAPE          & 4.60\% & 2.20\% \\
RMSE          & 25.58\,nm & 15.96\,nm \\
\bottomrule
\end{tabular}
\end{table}

The training curves (saved as \texttt{MSE\_images/mae.png}, \texttt{mape.png}, and \texttt{rmse.png}) show consistent convergence with no divergence between train and validation, confirming stable training dynamics.

%=====================================================================
\section{Inference Results}
%=====================================================================

Predictions were generated by the script \texttt{Predictions.py} (in \texttt{6\_results\_colab}) and saved as NumPy arrays under \texttt{MSE\_data/}. The evaluation script \texttt{1\_figures\_MSE.py} loads these arrays, computes error metrics (MAE, MSE, RMSE, MAPE, $R^2$), and performs a full statistical analysis of the error distributions, including:

\begin{itemize}
  \item Basic statistics (mean, median, standard deviation, variance, min, max, range, IQR)
  \item Distribution shape (skewness \& excess kurtosis with interpretation)
  \item Normality tests: Shapiro-Wilk (when $N \le 5000$), Kolmogorov-Smirnov, Anderson-Darling, Jarque-Bera
  \item Quantile comparison (empirical vs theoretical normal)
  \item Visual diagnostic plots (histogram with normal PDF overlaid, Q-Q plot, box plot, ECDF, moments bar chart, summary table)
  \item Outlier count (IQR method), entropy, coefficient of variation
\end{itemize}

All generated figures are stored under \texttt{MSE\_images/} in subdirectories \texttt{train/}, \texttt{val/}, \texttt{test/}, and \texttt{145/}.

%---------------------------------------------------------------------
\subsection{Train Split ($N = 960,000$)}
%---------------------------------------------------------------------

\begin{table}[H]
\centering
\caption{Error metrics — Train split.}
\label{tab:metrics_train}
\begin{tabular}{lc}
\toprule
\textbf{Metric} & \textbf{Value} \\
\midrule
MAE                 & 2.91\,nm \\
MSE                 & 14.73\,nm\textsuperscript{2} \\
RMSE                & 3.84\,nm \\
MAPE                & 0.61\% \\
$R^2$               & 0.9999 \\
\bottomrule
\end{tabular}
\end{table}

\subsubsection{APE Distribution}

\begin{itemize}
  \item \textbf{Mean APE}: 0.6120\%, \textbf{Median}: 0.3273\%, \textbf{Std}: 0.8593\%
  \item \textbf{Min}: 0.00\%, \textbf{Max}: 40.17\%, \textbf{IQR}: 0.5931\%
  \item \textbf{Skewness}: 6.2475 (highly skewed, right/positive)
  \item \textbf{Excess Kurtosis}: 105.3445 (leptokurtic, heavy-tailed)
  \item \textbf{Normality}: Rejected by all tests (KS $D = 0.2382$, $p \approx 0$; AD $A^2 = 90{,}671.85$; JB $= 4.50 \times 10^{8}$, $p \approx 0$)
  \item \textbf{Outliers} (IQR): 81,990 (8.54\%)
\end{itemize}

\subsubsection{PE Distribution}

\begin{itemize}
  \item \textbf{Mean PE}: $-0.3221$\%, \textbf{Median}: $-0.0481$\%, \textbf{Std}: 1.0047\%
  \item \textbf{Skewness}: $-3.6440$ (highly skewed, left/negative)
  \item \textbf{Excess Kurtosis}: 60.9873 (leptokurtic)
  \item \textbf{Normality}: Rejected by all tests
\end{itemize}

The negative mean PE ($-0.32\%$) indicates a negligible tendency to slightly over-predict thickness, concentrated on the left tail. However, the median ($-0.05\%$) is very close to zero, confirming that the bulk of predictions is well-centred.

%---------------------------------------------------------------------
\subsection{Validation Split ($N = 120,000$)}
%---------------------------------------------------------------------

\begin{table}[H]
\centering
\caption{Error metrics — Validation split.}
\label{tab:metrics_val}
\begin{tabular}{lc}
\toprule
\textbf{Metric} & \textbf{Value} \\
\midrule
MAE                 & 2.91\,nm \\
MSE                 & 14.80\,nm\textsuperscript{2} \\
RMSE                & 3.85\,nm \\
MAPE                & 0.61\% \\
$R^2$               & 0.9999 \\
\bottomrule
\end{tabular}
\end{table}

\subsubsection{APE Distribution}

\begin{itemize}
  \item \textbf{Mean APE}: 0.6095\%, \textbf{Median}: 0.3275\%, \textbf{Std}: 0.8552\%
  \item \textbf{Skewness}: 6.5218 (highly skewed, right)
  \item \textbf{Excess Kurtosis}: 124.7084 (leptokurtic)
  \item \textbf{Outliers}: 10,465 (8.72\%)
\end{itemize}

\subsubsection{PE Distribution}

\begin{itemize}
  \item \textbf{Mean PE}: $-0.3186$\%, \textbf{Median}: $-0.0446$\%
  \item \textbf{Skewness}: $-4.1251$ (highly skewed, left)
  \item \textbf{Excess Kurtosis}: 70.9283
\end{itemize}

%---------------------------------------------------------------------
\subsection{Test Split ($N = 120,000$)}
%---------------------------------------------------------------------

\begin{table}[H]
\centering
\caption{Error metrics — Test split.}
\label{tab:metrics_test}
\begin{tabular}{lc}
\toprule
\textbf{Metric} & \textbf{Value} \\
\midrule
MAE                 & 2.93\,nm \\
MSE                 & 15.14\,nm\textsuperscript{2} \\
RMSE                & 3.89\,nm \\
MAPE                & 0.62\% \\
$R^2$               & 0.9999 \\
\bottomrule
\end{tabular}
\end{table}

\subsubsection{APE Distribution}

\begin{itemize}
  \item \textbf{Mean APE}: 0.6153\%, \textbf{Median}: 0.3280\%, \textbf{Std}: 0.8716\%
  \item \textbf{Max APE}: 44.97\%
  \item \textbf{Skewness}: 7.2378 (highest among simulated splits)
  \item \textbf{Excess Kurtosis}: 167.2586 (heaviest-tailed among simulated splits)
  \item \textbf{Outliers}: 10,243 (8.54\%)
\end{itemize}

\subsubsection{PE Distribution}

\begin{itemize}
  \item \textbf{Mean PE}: $-0.3210$\%, \textbf{Median}: $-0.0468$\%
  \item \textbf{Skewness}: $-4.2738$, \textbf{Excess Kurtosis}: 95.4401
\end{itemize}

%---------------------------------------------------------------------
\subsection{Experimental 145\,nm ($N = 145$)}
%---------------------------------------------------------------------

\begin{table}[H]
\centering
\caption{Error metrics — Experimental 145\,nm.}
\label{tab:metrics_145}
\begin{tabular}{lc}
\toprule
\textbf{Metric} & \textbf{Value} \\
\midrule
MAE                 & 58.77\,nm \\
MSE                 & 7,029.82\,nm\textsuperscript{2} \\
RMSE                & 83.84\,nm \\
MAPE                & 7.38\% \\
$R^2$               & 0.9385 \\
\bottomrule
\end{tabular}
\end{table}

\subsubsection{APE Distribution}

\begin{itemize}
  \item \textbf{Mean APE}: 7.3797\%, \textbf{Median}: 6.4231\%, \textbf{Std}: 5.8508\%
  \item \textbf{Min}: 0.0049\%, \textbf{Max}: 27.6275\%
  \item \textbf{Skewness}: 0.8061 (moderately skewed, right)
  \item \textbf{Excess Kurtosis}: 0.1797 (mesokurtic — close to normal)
  \item \textbf{Normality}: Shapiro-Wilk $W = 0.9286$, $p = 1 \times 10^{-6}$ (rejected); KS $p = 0.044$ (borderline); Jarque-Bera $p = 0.0004$ (rejected)
  \item \textbf{Outliers}: 1 (0.69\%)
  \item \textbf{Coefficient of variation}: 0.7928
\end{itemize}

\subsubsection{PE Distribution}

\begin{itemize}
  \item \textbf{Mean PE}: $-1.7618$\%, \textbf{Median}: $-1.0253$\%, \textbf{Std}: 9.2706\%
  \item \textbf{Skewness}: $-0.0321$ (approximately symmetric)
  \item \textbf{Excess Kurtosis}: $-0.1377$ (mesokurtic)
  \item \textbf{Normality}: \textbf{Not rejected} by all four tests:
  \begin{itemize}
    \item Shapiro-Wilk: $W = 0.9941$, $p = 0.816$ \hfill (cannot reject normality)
    \item Kolmogorov-Smirnov: $D = 0.0720$, $p = 0.421$ \hfill (cannot reject)
    \item Anderson-Darling: $A^2 = 0.4087 < 0.7670$ \hfill (cannot reject at 5\%)
    \item Jarque-Bera: JB $= 0.2074$, $p = 0.902$ \hfill (cannot reject)
  \end{itemize}
  \item \textbf{Outliers}: 0 (0.00\%)
\end{itemize}

The fact that PE on the experimental set is approximately normally distributed is noteworthy: it suggests that the model's errors on real-world data, while larger in magnitude, behave as unbiased Gaussian noise. The negative mean PE ($-1.76\%$) indicates a small systematic under-prediction tendency for the 145\,nm experimental samples.

%=====================================================================
\section{Cross-Split Comparison}
%=====================================================================

Table~\ref{tab:comparison} summarises the key metrics across all four evaluation splits.

\begin{table}[H]
\centering
\caption{Cross-split error metrics comparison.}
\label{tab:comparison}
\begin{tabular}{lcccccc}
\toprule
\textbf{Split} & \textbf{$N$} & \textbf{MAE (nm)} & \textbf{RMSE (nm)} & \textbf{MAPE (\%)} & \textbf{$R^2$} & \textbf{Mean APE (\%)} \\
\midrule
Train      & 960,000 & 2.91 & 3.84 & 0.61 & 0.9999 & 0.6120 \\
Validation & 120,000 & 2.91 & 3.85 & 0.61 & 0.9999 & 0.6095 \\
Test       & 120,000 & 2.93 & 3.89 & 0.62 & 0.9999 & 0.6153 \\
145\,nm    &     145 & 58.77 & 83.84 & 7.38 & 0.9385 & 7.3797 \\
\bottomrule
\end{tabular}
\end{table}

\begin{itemize}
  \item \textbf{Simulated splits} (Train/Val/Test) are virtually indistinguishable across all metrics, providing strong evidence of robust generalization and absence of overfitting.
  \item \textbf{Experimental 145\,nm} shows a roughly 12$\times$ higher MAPE, consistent with the additional challenges of real-world spectra (noise, measurement uncertainty, model mismatch between simulation and experiment).
  \item The $R^2$ of 0.9385 on experimental data remains high, demonstrating that the model captures meaningful thickness variation even outside its training distribution.
\end{itemize}

%=====================================================================
\section{Distribution Analysis: Summary}
%=====================================================================

For all simulated splits, the APE (absolute percentage error) distributions share a characteristic shape:

\begin{itemize}
  \item Strongly positive-skewed (skewness 6.2–7.2)
  \item Extremely leptokurtic (excess kurtosis 105–167)
  \item Median APE is roughly half the mean APE ($\approx 0.33\%$ vs $\approx 0.61\%$), reflecting the heavy right tail
  \item Approximately 8.5\% of samples fall beyond the IQR-based outlier threshold
  \item Reject normality decisively across all tests ($p \approx 0$ for KS, AD, JB)
\end{itemize}

The PE distributions in simulated splits mirror this pattern on the left tail (negative skewness of $-3.6$ to $-4.3$, excess kurtosis 60–95), indicating that large errors (where they occur) tend to be over-predictions.

In contrast, the experimental 145\,nm set exhibits near-normal PE (Shapiro-Wilk $p = 0.816$), mild positive skew in APE, and mesokurtic behaviour in both PE and APE (excess kurtosis near zero). This shift from highly leptokurtic (simulated) to mesokurtic (experimental) may reflect the limited sample size ($N = 145$) or a fundamentally different error structure when applied to real data.

%=====================================================================
\section{Reproducibility}
%=====================================================================

\begin{itemize}
  \item \textbf{Model}: Pre-trained UNET saved as \texttt{.keras} file
  \item \textbf{Data}: Simulated spectra (train/val/test in \texttt{.h5} format) + experimental spectra from \texttt{dataframe\_spectrum\_thickness\_145\_final.pkl}
  \item \textbf{Training}: 100 epochs, Adam optimizer, learning rate $10^{-3}$, no learning rate schedule
  \item \textbf{Inference pipeline}:
  \begin{enumerate}
    \item \texttt{Predictions.py} (in \texttt{6\_results\_colab/}) generates \texttt{.npy} prediction files
    \item \texttt{training\_plots.py} generates training curve figures
    \item \texttt{1\_figures\_MSE.py} runs the full evaluation, generating all figures and \texttt{statistics.txt}
  \end{enumerate}
  \item \textbf{Output directory}: \texttt{MSE\_images/}
\end{itemize}

All scripts are designed to run in Google Colab with GPU acceleration. The paths reference Google Drive mounted at \texttt{/content/drive/}.

\end{document}
